\chapter{Tennis elbow}
\index{Tennis elbow}

\medskip
\noindent\textit{Educational reference; always seek professional advice for individual care.}

\section*{Definition and Overview}
Tennis elbow, known medically as lateral epicondylitis or lateral epicondylopathy, is a common and painful condition of the outer side of the elbow. It occurs where the tendons of the muscles that extend the wrist and fingers attach to the bony prominence on the outer side of the upper arm bone, the lateral epicondyle. The name is misleading: only a minority of people with the condition play tennis. Far more commonly it follows repetitive or heavy use of the forearm and wrist in everyday work and life, such as using hand tools, gardening, painting, typing, or repeatedly lifting with the palm facing down. The problem is essentially an overuse injury at the attachment point of the tendons, where repeated strain causes micro-tears, degeneration, and pain. The classic complaint is pain and tenderness over the outer elbow that is worsened by gripping, lifting, or twisting the wrist, and often a weakness that makes simple tasks like turning a key, holding a cup, or shaking hands uncomfortable. Tennis elbow affects adults of all ages, peaks between forty and sixty, and is equally common in men and women. The good news is that most cases improve substantially with conservative care, though the condition is notorious for taking months, sometimes up to a year or more, to settle fully. This chapter explains the causes, symptoms, diagnosis, and the evidence-based treatments that can speed recovery.

\section*{Epidemiology}
Tennis elbow is one of the commonest causes of elbow pain seen in general practice and one of the most frequent work-related musculoskeletal problems. It most often develops between the ages of forty and sixty, affects men and women about equally, and is rare under the age of thirty, although young, active people can develop it with heavy training. It is the dominant arm that is almost always affected, reflecting its role in most tasks. The strongest risk factors are repetitive wrist and forearm movements, forceful gripping, and heavy manual work: plumbers, painters, cooks, meatworkers, builders, and people who use vibrating tools or work with their hands for long periods are all at increased risk. Among sports, it is associated with racquet sports, particularly when technique or equipment is poor, and with golf and other gripping sports. The burden is considerable, because the pain interferes with both work and everyday activities, and the condition runs a long, fluctuating course that can test anyone's patience.

\section*{Aetiology and Pathophysiology}
Tennis elbow develops at the point where the common extensor tendon attaches to the lateral epicondyle of the humerus. This tendon is the shared attachment of the muscles that straighten (extend) the wrist and fingers, and it is loaded with every firm grip and every lift, because the wrist muscles brace the hand against the pull of the forearm flexors.

\begin{itemize}
    \item \textbf{Repetitive wrist extension:} repeated straightening of the wrist against resistance, as in racquet sports, hammering, sweeping, or using hand tools, strains the tendon attachment
    \item \textbf{Forceful gripping:} sustained or forceful gripping, such as carrying heavy objects, working with pliers, or using power tools, loads the extensor muscles and their tendon
    \item \textbf{Unaccustomed activity:} a sudden increase in gardening, DIY, sport, or work tasks overloads a tendon that is not conditioned for it
    \item \textbf{Poor technique and equipment:} a racquet with the wrong grip size or string tension, a poor backhand technique, or a poorly designed work tool increases the load on the elbow
    \item \textbf{Weakness and deconditioning:} weak wrist and forearm muscles leave the tendon to take the strain alone
    \item \textbf{Age-related tendon change:} the tendon's ability to repair declines with age, which is why the condition clusters in the over-forty age group
\end{itemize}

The tissue at the attachment becomes painful through a combination of micro-tears, degeneration of the collagen, and the growth of disorganised scar tissue, along with some inflammation, particularly early on. The condition is thus best understood as a tendinopathy at the elbow, and the treatment that works best follows the same principles as tendinopathy elsewhere: unload the tendon, then progressively load it again.

\section*{Clinical Features}
\begin{itemize}
    \item \textbf{Pain over the outer elbow:} aching or burning pain precisely over the bony lateral epicondyle, sometimes spreading down the forearm
    \item \textbf{Pain with gripping:} holding a cup, turning a doorknob or key, shaking hands, or lifting objects reproduces the pain
    \item \textbf{Pain with wrist movement:} straightening the wrist against resistance, or bending the wrist down while the elbow is straight, brings on the pain
    \item \textbf{Tenderness:} pressing directly on the lateral epicondyle is painful
    \item \textbf{Grip weakness:} many people feel their grip is weaker or painful, and they may drop things
    \item \textbf{Morning stiffness:} stiffness in the elbow and forearm on waking, easing with movement
    \item \textbf{Gradual onset:} the pain usually builds over weeks rather than following a single injury
    \item \textbf{Worsening with use:} tasks that involve lifting, twisting, or sustained gripping aggravate it, and rest eases it
\end{itemize}

In most cases the elbow is otherwise normal, with full movement and no swelling, redness, or numbness, which helps distinguish tennis elbow from other elbow problems.

\section*{Differential Diagnosis}
\begin{itemize}
    \item \textbf{Golfer's elbow (medial epicondylitis):} pain on the inner side of the elbow, at the attachment of the wrist flexors, caused by similar overuse in gripping and twisting with the palm down
    \item \textbf{Arthritis of the elbow joint:} osteoarthritis or inflammatory arthritis causes deeper joint pain, stiffness, and reduced movement rather than point tenderness at the epicondyle
    \item \textbf{Cervical spine problems:} pain referred from the neck can be felt in the elbow and forearm, often with neck stiffness or numbness
    \item \textbf{Nerve entrapment:} compression of the radial nerve, or of the ulnar nerve at the inner elbow, causes numbness, tingling, or weakness that tennis elbow does not
    \item \textbf{Biceps or triceps tendonitis:} pain at the front or back of the elbow with bending or straightening the arm
    \item \textbf{Elbow bursitis:} a swollen, boggy lump at the point of the elbow from inflammation of the olecranon bursa
    \item \textbf{Fracture or loose body:} a recent injury with bruising, or a joint that locks, may indicate a fracture or a loose piece of bone or cartilage
\end{itemize}

\section*{When to Seek Medical Care}
Most tennis elbow is managed at home and in primary care, but there are clear times to seek help.

Seek medical advice if you have:
\begin{itemize}
    \item Elbow pain that has lasted more than a few weeks despite rest and self-care
    \item Pain that is interfering with work, sport, or daily activities such as carrying, lifting, or holding
    \item Weakness or numbness in the hand or fingers
    \item Swelling, redness, or a lump at the elbow
    \item Pain after a specific injury or fall
\end{itemize}

\textbf{Contact your local emergency services or seek urgent care for:}
\begin{itemize}
    \item Severe pain with rapidly spreading redness, swelling, and fever, which can indicate infection
    \item Sudden severe weakness or deformity of the elbow after an injury, or an inability to use the arm
\end{itemize}

\section*{Diagnosis and Investigations}
Tennis elbow is a clinical diagnosis, and tests are only needed when the presentation is unusual, the problem is severe, or recovery is slow.

\begin{itemize}
    \item \textbf{History:} the location and behaviour of the pain, the activities that provoke it, work and sport demands, and how long it has been present
    \item \textbf{Examination:} pressing on the lateral epicondyle to confirm the tenderness, and reproducing the pain with resisted wrist extension and a resisted middle-finger test
    \item \textbf{Grip-strength testing:} an objective measure of the functional weakness and a useful way to monitor progress
    \item \textbf{Ultrasound:} shows the state of the tendon attachment, including thickening, tears, and calcification, and is used when a tear is suspected or an injection is planned
    \item \textbf{MRI:} reserved for persistent cases or when the diagnosis is unclear and other structures need to be examined
    \item \textbf{X-ray:} used after an injury to exclude a fracture, or to detect calcium deposits in a chronic case
    \item \textbf{Nerve conduction studies:} only if nerve entrapment is suspected on the basis of numbness or weakness
\end{itemize}

\section*{Management}
Most cases of tennis elbow settle with conservative treatment, and surgery is rarely needed. The core of treatment is unloading the tendon and then progressively rebuilding its strength.

\begin{itemize}
    \item \textbf{Activity modification:} reduce or modify the aggravating activities, including lifting, gripping, and repetitive wrist use; a band or brace just below the elbow can help off-load the tendon
    \item \textbf{Pain relief:} paracetamol, or a short course of an anti-inflammatory medicine such as ibuprofen, can help with pain and any early inflammation
    \item \textbf{Ice:} a cold pack over the tender point for fifteen to twenty minutes, especially after activity, helps settle pain
    \item \textbf{Gentle stretching:} stretching the wrist and forearm muscles, with the elbow straight and the wrist bent, keeps them flexible
    \item \textbf{Eccentric strengthening:} slowly lowering the wrist under a light weight, guided by a physiotherapist, is a core and effective exercise for tennis elbow
    \item \textbf{Progressive strengthening:} a graded program of wrist, forearm, and grip strengthening exercises, ideally under physiotherapy guidance, rebuilds the tendon
    \item \textbf{Physiotherapy:} can provide exercises, manual therapy, strapping, and advice on technique and ergonomics
    \item \textbf{Shockwave therapy:} has evidence of benefit for some people with chronic tennis elbow that has not responded to exercise
    \item \textbf{Corticosteroid injection:} gives good short-term relief but is followed by a high rate of recurrence, so it is used sparingly and with an exercise program
    \item \textbf{Surgery:} reserved for severe, long-standing cases that have failed a full course of non-surgical treatment; outcomes are generally good
\end{itemize}

\section*{Complications}
\begin{itemize}
    \item \textbf{Prolonged pain and disability:} the condition commonly lasts months, and some people are affected for a year or more, limiting work and leisure
    \item \textbf{Recurrence:} the problem returns readily if the cause of the overload is not addressed or rehabilitation is incomplete
    \item \textbf{Grip weakness:} ongoing pain and weakness can persist and affect daily tasks such as carrying groceries or holding a saucepan
    \item \textbf{Compensatory strain:} altered arm use can strain the shoulder, neck, or other elbow tendons
    \item \textbf{Treatment-related problems:} repeated corticosteroid injections can weaken the tendon and increase the chance of recurrence
    \item \textbf{Chronic tendon changes:} if the condition becomes chronic, the tendon attachment can develop degeneration and calcification that takes longer to treat
\end{itemize}

\section*{Prognosis and Outcomes}
Tennis elbow has a good long-term outlook, but it is slow. Most people improve within three to six months of consistent conservative treatment, though a substantial minority take a year or more to become fully pain-free, and the condition is well known for fluctuating: good weeks and bad weeks are normal. The evidence is clear that a structured exercise program gives the best results and that relying on rest alone or repeated injections tends to leave people with recurrences. In the long run, the majority of people recover completely with non-surgical care, and surgery is successful in most of the small group who need it. The keys to a good outcome are patience, a graduated strengthening program, and correcting the habits of work, sport, and equipment use that overloaded the tendon in the first place. For most people, the condition resolves, and they return to all their usual activities.

\section*{Prevention and Self-Care}
\begin{itemize}
    \item Warm up the forearm muscles and stretch before and after activities that load the wrist and elbow
    \item Build up the amount and force of gripping, lifting, and wrist work gradually rather than in a single burst
    \item Use correct technique and appropriate equipment, including a properly fitted racquet, and consider a counterforce brace during heavy activity
    \item Vary tasks and take regular breaks during long periods of repetitive manual work
    \item Strengthen the forearm and grip muscles as part of general fitness, and correct muscle imbalance with professional guidance
    \item Act on the first sign of elbow pain: modify the activity and start gentle stretching and strengthening rather than pushing through
    \item Use ergonomic tools, such as padded handles and reduced-vibration equipment, for work and gardening
\end{itemize}

\section*{Sources and Further Reading}
The following organisations provide reliable information about tennis elbow.

\begin{itemize}
    \item Sports Medicine Australia. \textit{Tennis elbow: information and rehabilitation.} https://sma.org.au/
    \item healthdirect Australia. \textit{Tennis elbow.} https://www.healthdirect.gov.au/
    \item Australian Physiotherapy Association. \textit{Lateral epicondylitis: a rehabilitation guide.} https://australian.physio/
    \item NHS. \textit{Tennis elbow.} https://www.nhs.uk/
    \item Mayo Clinic. \textit{Tennis elbow.} https://www.mayoclinic.org/
    \item Better Health Channel (Victoria). \textit{Tennis elbow.} https://www.betterhealth.vic.gov.au/
\end{itemize}

\clearpage
